\documentclass[11pt]{article}
% Shared local declarations. Document roots use the repository's business-admin family.
\usepackage{business-admin}
\usepackage{amsmath,amssymb,tabularx,longtable,array,booktabs,enumitem}
\usepackage{ragged2e}
\newcolumntype{Y}{>{\RaggedRight\arraybackslash}X}
\newcommand{\SuiteMeta}[2]{\def\DocTitle{#1}\def\DocSubtitle{#2}\def\DocVersion{1.0}\def\DocStatus{Educational and operational reference}\def\DocOwner{Parliamentary Procedure and Governance Suite}\def\DocSummary{Independently authored reference and working materials, based on RONR 12th edition.}}
\newcommand{\SuiteOpen}{\makedoctitle\vspace{-0.5em}\noindent\rule{\linewidth}{0.4pt}\par\medskip}
\newcommand{\RonrAnchor}[1]{\par\smallskip{\footnotesize\textbf{Primary reading:} Robert's Rules of Order Newly Revised, 12th ed., #1. Original synthesis; not official interpretation.}\par}
\newcommand{\WriteLine}{\par\noindent\rule{\linewidth}{0.35pt}\par\vspace{0.5cm}}
\newcommand{\WriteLines}[1]{\par\noindent\rule{\linewidth}{0.35pt}\par\vspace{0.35cm}\noindent\rule{\linewidth}{0.35pt}\par\vspace{0.35cm}\noindent\rule{\linewidth}{0.35pt}\par\vspace{0.25cm}}
\newcommand{\Field}[1]{\par\medskip\noindent\textbf{#1}\par\WriteLines{3}}
\newcommand{\GovCheckbox}{\(\square\)}
\setlength{\emergencystretch}{3em}

\SuiteMeta{Election Integrity Assessment — Template}{Evidence-led rubric | fictional assessment}
\begin{document}
\SuiteOpen
\section{Scoring methodology}
Use 0 = absent, 1 = inconsistent, 2 = partially implemented, 3 = implemented with evidence, 4 = independently verified and repeatable. Score only evidence within scope, note exclusions and cite governing authority. This maturity scale is original to the suite, not a published RONR maturity model.

\small\begin{longtable}{@{}p{.44\linewidth}cp{.39\linewidth}@{}}\toprule Criterion & Demo score & Fictional evidence\\\midrule\endhead

Elector eligibility verified & 2 & Draft control register; follow-up review required\\
Nomination method valid & 3 & Member portal entry and signed record\\
Ballot secrecy when required & 3 & Draft control register; follow-up review required\\
Ballots issued and returned reconciled & 2 & Member portal entry and signed record\\
Vote basis computed correctly & 3 & Draft control register; follow-up review required\\
Teller report independently reviewed & 3 & Member portal entry and signed record\\
No unauthorized plurality election & 2 & Draft control register; follow-up review required\\
Transition documented & 3 & Member portal entry and signed record\\
\bottomrule\end{longtable}\normalsize
\section{Illustrative rating}
Fictional demonstration score: 21 of 32. Treat any zero on a legal validity or due-process gate as a mandatory escalation, not as an offset by high scores elsewhere.

\section{Corrective-action protocol}
\begin{enumerate}[leftmargin=*,itemsep=4pt]
\item Attach source or evidence ID for each scored criterion.
\item Document weakness, responsible owner, due date and retest method.
\item Obtain an independent review for contested interpretation or disciplinary concerns.
\item Keep real sensitive evidence outside the public repository.
\end{enumerate}
\end{document}
